% Shared problem page; stochastic oracle on page 2, experiment results on page 3.
{\color{ink}\fontsize{19}{22}\selectfont\bfseries The Gromov--Wasserstein problem\par}
{\small\color{muted}Reference page \enspace / \enspace the optimization problem\par}
\vspace{8pt}

\textbf{Match relationships when features cannot be compared directly.}
Suppose two datasets describe different collections of objects, or two different
measurements of the same objects. Their feature spaces may be incompatible.
Gromov--Wasserstein (GW) matching uses the distances \emph{within} each dataset:
a good correspondence pairs objects so that their relationships are preserved.
It does not require a distance between an object in one feature space and an
object in the other.

\textbf{Data and unknown.}
Let $A\in\mathbb R^{n\times n}$ and $B\in\mathbb R^{m\times m}$ be the two
symmetric distance matrices. Positive probability vectors $a$ and $b$ specify
the mass of each object; uniform weights give $a_i=1/n$ and $b_j=1/m$.
The unknown is a \textbf{soft matching} $P\in\mathbb R^{n\times m}$:
$P_{ij}$ is the mass assigned between source object $i$ and target object $j$.
Balanced matching requires
\[
 U(a,b)=\{P\geq0:\ P\one_m=a,\ P^\top\one_n=b\}.
\]
The row and column sums prescribe how much each object contributes. They do not
force a one-to-one matching.

\textbf{Objective: preserve pairwise distances.}
The unregularized squared GW problem used here is
\[
 \boxed{\min_{P\in U(a,b)} f(P)},\qquad
 f(P)=\sum_{i,k=1}^{n}\sum_{j,\ell=1}^{m}
       (A_{ik}-B_{j\ell})^2 P_{ij}P_{k\ell}.
\]
Two simultaneous matches, $i\leftrightarrow j$ and $k\leftrightarrow\ell$,
incur a cost when $A_{ik}$ and $B_{j\ell}$ disagree. Equivalently, $f(P)$ is the
expected squared distance discrepancy from two independent draws from $P$.
This convention includes neither a factor $1/2$ nor entropy regularization.
The objective is smooth and generally nonconvex; the feasible transport
polytope is convex. A small loss alone does not certify a global optimum or a
scientifically correct correspondence.

\textbf{Composite form for FraMES.}
Separate the inexpensive row constraints from the column constraint:
\[
 \begin{aligned}
 C&=\{P\geq0:P\one_m=a\}, & T(P)&=P^\top\one_n,\\
 g(y)&=\iota_{\{b\}}(y)
 =\begin{cases}0,&y=b,\\+\infty,&y\ne b,\end{cases}
 &\min_{P\in C}\;&f(P)+g(TP).
 \end{aligned}
\]
\textbf{Smoothing the column constraint and its gradient.}
The indicator $g$ is not differentiable; the smooth quantity is its Moreau
envelope. Projection onto the singleton $\{b\}$ gives
\[
 \operatorname{prox}_{\mu g}(y)=b,\qquad
 g_\mu(y)=\frac{\|y-b\|_2^2}{2\mu},\qquad
 \nabla g_\mu(y)=\frac{y-b}{\mu}.
\]
Under Frobenius and Euclidean inner products, $T^*u=\one_nu^\top$.
The chain rule therefore gives the explicit penalty gradient
\[
 \boxed{\nabla(g_\mu\circ T)(P)
 =\frac{T^*(TP-b)}{\mu}
 =\frac{\one_n(P^\top\one_n-b)^\top}{\mu}.}
\]
This term is computed directly from the current column residual, without random
probes. A smaller $\mu$ strengthens the penalty for a fixed residual.

FraMES combines a momentum estimate of $\nabla f(P)$ with this current penalty
gradient. Its linear minimization oracle places each row's mass in a minimizing
column. Convex-combination updates preserve nonnegativity and row masses;
column masses can be violated during optimization.

\textbf{What a results page measures.}
Feasible loss is the full $f$ evaluated after output-only transport repair to
numerical tolerance. Raw loss, feasibility error and the smoothed penalty are
separate quantities. Loss magnitudes depend on distance scaling, so reductions
are interpreted within a fixed dataset. The next page explains the stochastic
oracle for $\nabla f$.
\vfill
{\footnotesize\color{muted}Shared background. Page 2 explains the stochastic oracle; page 3 describes the experiment's data, settings and results.}
